\documentclass[10pt,a4paper]{amsart}
\usepackage[margin=19mm]{geometry}
\usepackage{amsmath,amssymb,mathtools}
\usepackage[scr=rsfs,cal=euler]{mathalfa}
\usepackage{xfrac}
\usepackage[unicode,hidelinks,bookmarks=false]{hyperref}
\newcommand{\ie}{i.e.\ }
\newcommand{\cf}{cf.\ }
\newcommand{\resp}{respectively\ }
\newcommand{\Subsection}{\subsection}
\providecommand{\nobreakdash}{\nobreak\hbox{-}}
\setlength{\emergencystretch}{2em}
\multlinegap0pt
% ---------------------------------------------------------------------------
% Editorial AMS wrapper prelude. The theorem environments and the mathematical macros below are
% transcribed from the paper's own preamble (cached-originals/source0.tex lines 170--174); the AMS amsart
% class, the named format packages of the pinned TeX Live runtime and this editorial formatting are not
% part of the author's text. No mathematics is changed.
% ---------------------------------------------------------------------------
\usepackage{stmaryrd}
\newtheorem{theorem}{Theorem}[section]
\newtheorem{lemma}[theorem]{Lemma}
\newtheorem{proposition}[theorem]{Proposition}
\newtheorem{definition}[theorem]{Definition}
\newtheorem{remark}[theorem]{Remark}

\begin{document}
\title{Explosivity in 1-d Activated Random Walk}
\author{Nicolas Forien, Christopher Hoffman, Tobias Johnson, Josh Meisel, Jacob Richey and Leonardo T. Rolla}
\date{Source: arXiv:2601.03411v2; distributed under CC BY 4.0}
\maketitle
\noindent\emph{Source, version and licence.} \emph{Explosivity in 1-d Activated Random Walk}, by Nicolas Forien, Christopher Hoffman, Tobias Johnson, Josh Meisel, Jacob Richey and Leonardo T. Rolla. The source of this editable unit is the e-print of \textbf{version v2} of arXiv:2601.03411, https://arxiv.org/abs/2601.03411v2, whose material is distributed under the Creative Commons Attribution 4.0 International licence, http://creativecommons.org/licenses/by/4.0/, as recorded in the arXiv licence metadata for this paper and shown on that abstract page. The whole of the body below is version v2, the newest version of the paper. Full rights notice: licenses/arxiv-2601.03411-CC-BY-4.0.txt.\par\smallskip
\noindent\textit{Editorial scope.} Counted unit: original Theorem 1.1 of the article, the statement carried by the source environment \texttt{theorem} with source label \texttt{thm:explosive}, cached-originals/source0.tex lines 267--274, printed as \textbf{Theorem 1.1} exactly as in the article. It is delivered together with the authors' own complete proof as the single primary proof body, source0.tex lines 626--640 (15 lines): the \texttt{proof} environment whose optional argument names it "Proof of Theorem \textbackslash ref\{thm:explosive\}", which treats the subcritical and supercritical cases separately and closes with the dichotomy of the statement. No proof marker is added and no mathematics is written, repaired or re-derived; the authors' notation and the printed slips of the source are kept literal. Necessary complete context is retained uncounted, in source order: source0.tex lines 383--389, Lemma 2.1 on preemptive waking, which the supercritical half of the counted proof invokes; lines 418--436, Definition 2.2 of the odometer processes $X$ and $Y$; lines 439--465, Proposition 2.3 on right avalanches, together with its complete proof at lines 467--488; lines 495--517, Lemma 2.4 on averages, with its complete proof at lines 519--547; and lines 550--568, Proposition 2.5, the detection criterion that the counted proof cites, with its complete proof at lines 570--624. The article's own numbering is reproduced by explicit counter settings: the article's publisher class declares \texttt{\textbackslash newtheorem\{theorem\}\{Theorem\}[section]} with lemma, proposition, definition and remark sharing that counter, so the wrapper restates those declarations and sets the counter before each retained passage; the counted statement prints as Theorem 1.1, the preemptive-waking lemma as Lemma 2.1, the odometer definition as Definition 2.2 and the two propositions as Propositions 2.3 and 2.5, exactly as in the article. Every \texttt{\textbackslash ref} and \texttt{\textbackslash eqref} inside every retained passage resolves inside this delivered file, and the two citations that occur in the retained passages, \texttt{RollaSidoraviciusZindy19} and \texttt{rolla2020activated}, are delivered as the authors' own bibliography entries transcribed verbatim from the article's own \texttt{thebibliography}, keeping the authors' own printed numbers. The retained passages contain no figure environment and no graphics-inclusion command, so no artwork is redistributed. The source's own class is the publisher class \texttt{ejpecp} of the Electronic Journal of Probability (\texttt{\textbackslash documentclass[ECP]\{ejpecp\}}); that class file is shipped by the e-print but is \emph{not} shipped, loaded or used by this package and is never redistributed, and the wrapper is the AMS \texttt{amsart} class with the article's theorem declarations restated and no article macros (the article declares none). Every declared body of this unit is an exact, unmodified span of source0.tex: no body is a bounded passage comparison, \texttt{omit} was not used anywhere, and consequently no fidelity receipt is asserted in delivery-evidence.json. The complete concatenated original source is bundled as sources/192223/source0.tex. Omitted from this editable unit and therefore absent from it are source0.tex lines 1--266; 275--382; 390--417; 437--438; 466--494; 518--549; 569--569; 625--625; 641--722; 729--735; 742--767, that is the article's own preamble, author block and abstract, the introduction, the remaining sections and remarks, and the remaining bibliography entries, all of which remain present in the bundled source but are not delivered here. The AMS \texttt{amsart} wrapper, the named format packages of the pinned TeX Live runtime and this editorial source, licence and scope paragraph are editorial formatting only.\par\smallskip
\setcounter{section}{1}
\setcounter{equation}{0}
\setcounter{theorem}{0}
\begin{theorem}%
%%LEAP%%%\label{thm1.1}
\label{thm:explosive}
Let $\sigma $ be a random element in
$(\mathbb{N}\cup \{\mathfrak{s}\})^{\mathbb{Z}}$ with an ergodic distribution
and density $\mathbf{E}\mathopen{\lvert }\sigma (0)\mathclose{\rvert } = \rho $. Then $\sigma $ is a.s.\ explosive
if $\rho > \rho _{c}$ and a.s.\ not explosive if $\rho < \rho _{c}$.
\end{theorem}

\setcounter{section}{2}
\setcounter{equation}{0}
\setcounter{theorem}{0}
\begin{lemma}%
%%LEAP%%%\label{lem2.1}
\label{lem:preemptive}
Let $V\subseteq \mathbb{Z}$, and suppose that $U\subseteq V$ is a collection
of sites visited during the stabilization of $\sigma $ on $V$. Then stabilizing
$\sigma $ or $\mathsf{A}^{U}\sigma $ on $V$ yields the same odometer.
\end{lemma}

\setcounter{section}{2}
\setcounter{equation}{0}
\setcounter{theorem}{1}
\begin{definition}[$X(n, \sigma )$, $Y(n, \sigma )$, $X(n, \sigma ,u_{0})$, $Y(n,\sigma ,u_{0})$]
\label{defn2.2}
Given $n\ge 1$ and a particle configuration $\sigma $, we define
$n<X(n, \sigma )\leq \infty $ as the leftmost unvisited site in
$\llbracket n+1,\infty )$ when stabilizing
$\mathsf{A}^{\mathopen{\llbracket }1,n \mathclose{\rrbracket }}\sigma $ on $\mathbb{Z}_{>0}$. We let
$X(n,\sigma )=\infty $ when all sites in $\llbracket n+1,\infty )$ are
visited. Similarly, we define $-\infty \le Y(n, \sigma ) < -n$ as the rightmost
unvisited site in $(-\infty ,-n-1\rrbracket $ when stabilizing
$\mathsf{A}^{\mathopen{\llbracket }-n,-1\mathclose{\rrbracket }}\sigma $ on $\mathbb{Z}_{<0}$.

We will sometimes wish to consider these random variables starting midstream
in an ARW system, i.e., after some instructions have already been executed.
For a given odometer $u \colon \mathbb{Z}\to \mathbb{N}$, we write
$X(n, \sigma , u)$ and $Y(n, \sigma , u)$ to indicate the corresponding
quantities for the system in which the first instruction executed at site~$i$
begins at index $u(i)$ rather than $0$. The initial configuration is still
taken to be $\sigma $.
\end{definition}

\setcounter{section}{2}
\setcounter{equation}{0}
\setcounter{theorem}{2}
\begin{proposition}%
%%LEAP%%%\label{prop2.3}
\label{prop:right.avalanche}
For some $\epsilon ,\beta >0$ and some $N\geq 1$, let
$\sigma \colon \mathbb{Z}\to \mathbb{N}\cup \{\mathfrak{s}\}$ be a deterministic
configuration, which for all $n \ge N$ satisfies
%
%e2.1 #&#
\begin{align}
\sum _{j=1}^{n} j\mathopen{\lvert }\sigma (j)\mathclose{\rvert } &\geq
\frac{(\rho _{c}+\epsilon )n^{2}}{2}
\label{eq:center.of.mass}
%%LEAP%%%\label{eq2.1}
\\
\intertext{and} \sum _{j=1}^{n}\mathopen{\lvert }\sigma (j)\mathclose{\rvert }&\leq \beta n.
\label{eq:max.density}
\end{align}
%
Then for constants $c,C>0$ depending only on $\epsilon $, $\beta $ and
$\lambda $,
%
\begin{align*}
\mathbf{P}(X(n, \sigma )<\infty ) &\leq Ce^{-cn},
\end{align*}
%
for all $n \ge N$.
\end{proposition}

\setcounter{section}{2}
\setcounter{equation}{2}
\setcounter{theorem}{3}
\begin{lemma}%
%%LEAP%%%\label{lem2.4}
\label{lem:averages}
Let $a_{1},a_{2},\ldots \in \mathbb{R}$ and suppose that
%
%e2.3 #&#
\begin{align}
\label{eq:erg}
\lim _{n\to \infty}\frac{1}{n}\sum _{j=1}^{n}a_{j}=\rho .
%%LEAP%%%\label{eq2.3}
\end{align}
%
Then
%
%e2.4 #&#
\begin{align}
\label{eq:weighted}
\lim _{n\to \infty}\frac{1}{n^{2}}\sum _{j=1}^{n}ja_{j}=
\frac{\rho }{2}.
%%LEAP%%%\label{eq2.4}
\end{align}
%
\end{lemma}

\setcounter{section}{2}
\setcounter{equation}{4}
\setcounter{theorem}{4}
\begin{proposition}%
%%LEAP%%%\label{prop2.5}
\label{prop:det.explosive}
Let $\sigma \colon \mathbb{Z}\to \mathbb{N}\cup \{\mathfrak{s}\}$ be a
configuration of particles that contains at least one active particle and
satisfies
%
%e2.5 #&#
\begin{align}
\label{eq:det.explosive}
\lim _{n\to \infty}\frac {1} {n} \sum _{j=1}^{n}\mathopen{\lvert }\sigma (j)\mathclose{\rvert }=
\lim _{n\to \infty}\frac {1} {n} \sum _{j=-n}^{-1}\mathopen{\lvert }\sigma (j)\mathclose{\rvert }>
\rho _{c}.
%%LEAP%%%\label{eq2.5}
\end{align}
%
Then it occurs with positive probability that all sites in
$\mathbb{Z}$ are visited when starting ARW from $\sigma $.
\end{proposition}

\setcounter{section}{2}
\setcounter{equation}{5}
\setcounter{theorem}{5}
\begin{proof}
By shifting $\sigma $, we can take it to have an active particle at site~$0$
without loss of generality. Let $\rho $ be the value of the limits in
\eqref{eq:det.explosive}. By Lemma~\ref{lem:averages},
%
\begin{align*}
\lim _{n\to \infty}\frac{1}{n^{2}}\sum _{j=1}^{n} j\mathopen{\lvert }\sigma (j)\mathclose{\rvert }=
\frac{\rho }{2}.
\end{align*}
%
Thus, for $\epsilon =(\rho -\rho _{c})/2$ and $\beta = 2\rho $, there exists
some $N$ large enough that the conditions of Proposition~\ref{prop:right.avalanche}
hold, and we can conclude that
%
\begin{align*}
\mathbf{P}(X(n, \sigma )<\infty ) &\leq Ce^{-cn}
\end{align*}
%
for all $n\geq N$. By the same argument applied to $\sigma $ flipped across
the origin, there exists some $N'$ such that
%
\begin{align*}
\mathbf{P}(Y(n, \sigma )>-\infty ) &\leq Ce^{-cn}
\end{align*}
%
for all $n\geq N'$, with the same constants $c,C>0$.

Choose $m$ larger than $N$ and $N'$ and satisfying $Ce^{-cm}<1/2$. We stabilize
$\sigma $ on $\mathbb{Z}$ as follows. First, topple an active particle
at the origin until it falls asleep. Let $V_{0}\subseteq \mathbb{Z}$ be
the interval of sites it visits, and let $u_{0}$ be the system's odometer
after it sleeps. Conditional on $V_{0} \supseteq \mathopen{\llbracket }-m,m \mathclose{\rrbracket }$, the random
variables $X(m, \sigma , u_{0})$ and $Y(m, \sigma , u_{0})$ are distributed
as $X(m,\sigma )$ and $Y(m,\sigma )$, respectively, since starting the
instruction stacks at the indices given by $u_{0}$ does not alter their
distributions.

Now, the event $\{V_{0} \supseteq \mathopen{\llbracket }-m,m \mathclose{\rrbracket }\}$ holds with positive probability,
and
%
\begin{align*}
\mathbf{P}\bigl(X(m, \sigma , u_{0})<\infty \;\big \vert \;V_{0}
\supseteq \mathopen{\llbracket }-m,m \mathclose{\rrbracket }\bigr)&\leq Ce^{-cm}<1/2
\\
\intertext{and} \mathbf{P}\bigl(Y(m, \sigma , u_{0})>-\infty \;
\big \vert \;V_{0} \supseteq \mathopen{\llbracket }-m,m \mathclose{\rrbracket }\bigr)&\leq Ce^{-cm}<1/2.
\end{align*}
%
Thus by a union bound, it occurs with positive probability conditional
on the event $\{V_{0} \supseteq \mathopen{\llbracket }-m,m \mathclose{\rrbracket }\}$ that events
$\{X(m, \sigma , u_{0})=\infty \}$ and
$\{Y(m, \sigma , u_{0})=-\infty \}$ both occur. When all three events occur,
which occurs with positive probability, all sites in $\mathbb{Z}$ are visited
during the stabilization of $\sigma $, by definition of $X$ and $Y$.
\end{proof}

\setcounter{section}{2}
\setcounter{equation}{5}
\setcounter{theorem}{5}
\begin{proof}[Proof of Theorem~\ref{thm:explosive}]
If $\rho <\rho _{c}$, then ARW fixates even starting from $\sigma $ with
all particles activated from the beginning by
\cite[Theorem~1]{RollaSidoraviciusZindy19}. Monotonicity of the odometer
when activating particles \cite[Lemma~2.5]{rolla2020activated} implies
that $\sigma $ is not explosive in this case.

If $\rho >\rho _{c}$, then by the ergodic theorem and Proposition~\ref{prop:det.explosive},
stabilizing the configuration $\sigma $ after activating the particles
at any nonvacant site visits all of $\mathbb{Z}$ with positive probability.
On this event, the odometer stabilizing $\sigma$ on $\mathbb{Z}$ is unchanged
by preemptively waking all particles in $\sigma$ by Lemma~\ref{lem:preemptive}.
But by \cite[Theorem~1]{RollaSidoraviciusZindy19}, that odometer is infinite
a.s., and so fixation does not occur.
\end{proof}

\begin{thebibliography}{99}
\bibitem[10]{rolla2020activated} Leonardo~T. Rolla, \emph{Activated random walks on {$\mathbb{Z}^{d}$}}, Probab. Surv. \textbf{17} (2020), 478--544. \MR{4152668} %b11 #&#
\bibitem[12]{RollaSidoraviciusZindy19} Leonardo~T. Rolla, Vladas Sidoravicius, and Olivier Zindy, \emph{Universality and sharpness in activated random walks}, Ann. Henri Poincar{\'e} \textbf{20} (2019), 1823--1835. %b13 #&#
\end{thebibliography}
\end{document}
